\documentclass[a4paper,12pt]{article}
\usepackage[T2A]{fontenc}
\usepackage[utf8]{inputenc}
\usepackage[russian]{babel}
\usepackage{amsmath,amssymb,mathtools}
\usepackage{PTSans}
\renewcommand{\familydefault}{\sfdefault}
\usepackage{icomma}
\usepackage{geometry}
\geometry{top=18mm,bottom=18mm,left=18mm,right=18mm}
\usepackage{microtype}
\usepackage{tikz}
\usetikzlibrary{calc}
\usepackage{tabularx,booktabs,array}
\usepackage{enumitem}
\usepackage{multicol}
\usepackage{parskip}
\usepackage{fancyhdr}
\usepackage{setspace}
\usepackage{xcolor}
\usepackage{titlesec}

\definecolor{accent}{HTML}{1F6F78}
\definecolor{darkaccent}{HTML}{174F56}
\definecolor{textgray}{HTML}{2E3335}
\definecolor{softgray}{HTML}{EEF3F4}
\definecolor{warn}{HTML}{9A5B34}

\setlength{\parindent}{0pt}
\setlength{\parskip}{0.55em}
\setlength{\headheight}{25pt}
\onehalfspacing
\color{textgray}
\setlist[itemize]{leftmargin=1.35em,itemsep=0.22em,topsep=0.25em}
\setlist[enumerate]{leftmargin=1.55em,itemsep=0.28em,topsep=0.3em}
\titleformat{\section}{\Large\bfseries\color{darkaccent}}{}{0pt}{}[\vspace{-0.2em}\color{accent}\rule{\linewidth}{0.6pt}]
\titleformat{\subsection}{\large\bfseries\color{darkaccent}}{}{0pt}{}
\titlespacing*{\section}{0pt}{1.0em}{0.6em}
\titlespacing*{\subsection}{0pt}{0.8em}{0.25em}

\pagestyle{fancy}
\fancyhf{}
\fancyhead[L]{\small\color{darkaccent}Математика · чек-лист}
\fancyhead[R]{\small\color{darkaccent}03.10.26}
\fancyfoot[C]{\small\color{textgray}\thepage}
\renewcommand{\headrulewidth}{0.4pt}
\renewcommand{\headrule}{\hbox to\headwidth{\color{accent}\leaders\hrule height \headrulewidth\hfill}}

\newcommand{\checkitem}[1]{\item[$\square$] #1}
\newcommand{\nodtext}{\text{НОД}}
\newcommand{\noktext}{\text{НОК}}
\newcommand{\key}[1]{\textbf{\color{darkaccent}#1}}
\newcommand{\important}[1]{%
  \par\vspace{0.25em}
  \noindent\colorbox{softgray}{\parbox{0.965\linewidth}{\textbf{#1}}}\par\vspace{0.25em}}
\newcommand{\errornote}[1]{%
  \par\vspace{0.2em}
  \noindent\textcolor{warn}{\textbf{Типичная ошибка.}} #1\par}

\begin{document}

\begin{titlepage}
\thispagestyle{empty}
\centering
\vspace*{22mm}

{\large\color{accent}\textbf{МАТЕМАТИКА}}\\[10mm]
{\fontsize{30}{34}\selectfont\bfseries\color{darkaccent}Чек-лист}\\[4mm]
{\LARGE\bfseries Делимость натуральных чисел, НОД и НОК}\\[5mm]
{\large Делители · простые множители · чётность · дроби}\\[16mm]

\begin{minipage}{0.78\textwidth}
\centering
\large
Короткое повторение перед контрольной: определения, признаки, алгоритмы, показательные примеры и места, где чаще всего возникает ошибка.
\end{minipage}

\vfill
{\large 03.10.26}\\[8mm]
{\normalsize\bfseries Лёвин Артём Александрович}\\[2mm]
{\normalsize эксклюзивно для Нади Клименко}

\vspace{15mm}
\begin{tikzpicture}[remember picture,overlay]
  \draw[accent,line width=1.15pt]
    ($(current page.south west)+(14mm,18mm)$)
    .. controls ($(current page.south west)+(62mm,29mm)$)
    and ($(current page.south east)+(-70mm,8mm)$)
    .. ($(current page.south east)+(-14mm,21mm)$);
  \draw[accent!55,line width=0.65pt]
    ($(current page.south west)+(20mm,12mm)$)
    .. controls ($(current page.south west)+(76mm,22mm)$)
    and ($(current page.south east)+(-56mm,3mm)$)
    .. ($(current page.south east)+(-19mm,14mm)$);
\end{tikzpicture}
\end{titlepage}

\section{1. Делимость: что нужно видеть сразу}

\subsection{Делитель и кратное}
Если натуральное число $a$ делится на $b$ без остатка, то $b$ --- \key{делитель} числа $a$, а $a$ --- \key{кратное} числа $b$.

Например, $6\mid 36$, поэтому $6$ --- делитель $36$, а $36$ --- кратное $6$.

\important{У любого натурального числа есть делители $1$ и само это число. Число $1$ не является ни простым, ни составным.}

\subsection{Признаки делимости, которые нужны на контрольной}
\renewcommand{\arraystretch}{1.28}
\begin{tabularx}{\linewidth}{>{\bfseries\color{darkaccent}}p{0.12\linewidth} X}
\toprule
На & Признак \\
\midrule
$2$ & последняя цифра чётная: $0,2,4,6,8$; \\
$3$ & сумма цифр делится на $3$; \\
$4$ & число, образованное двумя последними цифрами, делится на $4$; \\
$5$ & последняя цифра $0$ или $5$; \\
$6$ & число одновременно делится на $2$ и на $3$; \\
$9$ & сумма цифр делится на $9$; \\
$12$ & число одновременно делится на $3$ и на $4$; \\
$18$ & число одновременно делится на $2$ и на $9$. \\
\bottomrule
\end{tabularx}

\errornote{Для делимости на $4$ проверяют \emph{число из двух последних цифр}, а не каждую из этих цифр отдельно. Например, $144$ делится на $4$, потому что $44$ делится на $4$.}

\subsection{Быстрая проверка на одном наборе}
Для чисел $144,\ 315,\ 726,\ 935,\ 1248,\ 1725$:
\begin{itemize}
  \item на $4$ делятся $144$ и $1248$;
  \item на $6$ делятся $144$, $726$, $1248$;
  \item на $9$ делятся $144$ и $315$;
  \item на $18$ из этого набора делится $144$.
\end{itemize}

\section{2. Делители и простые множители}

\subsection{Простые и составные числа}
\key{Простое число} имеет ровно два натуральных делителя: $1$ и само число.\\
\key{Составное число} имеет больше двух натуральных делителей.

Для числа $36$:
\[
1,\ 2,\ 3,\ 4,\ 6,\ 9,\ 12,\ 18,\ 36.
\]
Простые среди них: $2$ и $3$.

\important{При поиске делителей удобно искать их парами: $1\cdot36$, $2\cdot18$, $3\cdot12$, $4\cdot9$, $6\cdot6$. После $\sqrt{36}=6$ новые пары уже не появятся.}

\subsection{Разложение на простые множители}
Делим число последовательно на простые числа $2,3,5,7,\ldots$, пока не получим $1$.

Показательный пример:
\[
756=2\cdot378=2^2\cdot189=2^2\cdot3^3\cdot7.
\]
То есть
\[
\boxed{756=2^2\cdot3^3\cdot7}.
\]

\subsection{Количество натуральных делителей}
Если
\[
n=p_1^{\alpha_1}p_2^{\alpha_2}\cdots p_k^{\alpha_k},
\]
то количество натуральных делителей равно
\[
\boxed{\tau(n)=(\alpha_1+1)(\alpha_2+1)\cdots(\alpha_k+1)}.
\]

Для $756=2^2\cdot3^3\cdot7^1$:
\[
\tau(756)=(2+1)(3+1)(1+1)=3\cdot4\cdot2=24.
\]

Проверка на знакомом числе:
\[
36=2^2\cdot3^2\quad\Rightarrow\quad \tau(36)=(2+1)(2+1)=9.
\]

\errornote{В формуле прибавляют единицу к \emph{показателям степеней}, а не к самим простым множителям.}

\subsection{Взаимно простые числа}
Два числа взаимно просты, если их единственный общий натуральный делитель --- $1$, то есть
\[
\nodtext(a,b)=1.
\]
Например, $14$ и $25$ взаимно просты.

\section{3. НОД и НОК}

\subsection{Через разложение на простые множители}
Пусть числа разложены на простые множители.

\begin{itemize}
  \item \key{НОД}: берём только общие простые множители в \textbf{минимальных} степенях.
  \item \key{НОК}: берём все встречающиеся простые множители в \textbf{максимальных} степенях.
\end{itemize}

\important{Полезная проверка: $\nodtext(a,b)\cdot\noktext(a,b)=a\cdot b$.}

Например, поскольку $84\mid252$,
\[
\nodtext(252,84)=84,\qquad
\noktext(252,84)=252.
\]

\subsection{Алгоритм Евклида для больших чисел}
Когда разложение на простые множители становится неудобным, используйте остатки от деления:
\begin{enumerate}
  \item большее число делим на меньшее;
  \item вместо пары чисел берём \emph{меньшее число и остаток};
  \item повторяем, пока остаток не станет равен нулю;
  \item последний ненулевой остаток --- НОД.
\end{enumerate}

Пример:
\begin{align*}
1220&=2\cdot440+340,\\
440&=1\cdot340+100,\\
340&=3\cdot100+40,\\
100&=2\cdot40+20,\\
40&=2\cdot20+0.
\end{align*}
Следовательно,
\[
\boxed{\nodtext(1220,440)=20}.
\]

\errornote{Останавливаемся только после появления нулевого остатка. Ответом является предыдущий, последний ненулевой остаток.}

\subsection{Как выбирать НОД или НОК}
\begin{tabularx}{\linewidth}{>{\bfseries\color{darkaccent}}p{0.22\linewidth} X}
\toprule
Ищем НОД & Когда нужно максимально крупно \emph{разделить} объекты на одинаковые части, сократить дробь, найти общий делитель. \\
\midrule
Ищем НОК & Когда нужно найти первый общий момент/шаг/кратное, общий знаменатель или число, кратное нескольким данным числам. \\
\bottomrule
\end{tabularx}

\section{4. Чётность как быстрый инструмент}

Вместо вычислений часто достаточно определить чётность каждого слагаемого или множителя.

\begin{tabularx}{\linewidth}{>{\centering\arraybackslash}p{0.22\linewidth} >{\centering\arraybackslash}p{0.22\linewidth} X}
\toprule
Первое & Второе & Результат \\
\midrule
чётное & чётное & сумма чётная \\
нечётное & нечётное & сумма чётная \\
чётное & нечётное & сумма нечётная \\
чётное & любое целое & произведение чётное \\
нечётное & нечётное & произведение нечётное \\
\bottomrule
\end{tabularx}

Отсюда:
\begin{itemize}
  \item сумма чётного количества нечётных слагаемых --- чётная;
  \item сумма нечётного количества нечётных слагаемых --- нечётная;
  \item произведение чётно, если среди множителей есть хотя бы один чётный.
\end{itemize}

Например,
\[
15\cdot17+24\cdot30
\]
не нужно вычислять: первое произведение нечётное, второе чётное, поэтому сумма нечётная.

\section{5. НОД и НОК в действиях с дробями}

\subsection{Общий знаменатель через НОК}
Для сложения или вычитания дробей с разными знаменателями удобно брать \key{НОК знаменателей}.

Пример:
\[
\frac{1}{210}+\frac{3}{104}.
\]
Разложения:
\[
210=2\cdot3\cdot5\cdot7,\qquad 104=2^3\cdot13.
\]
Поэтому
\[
\noktext(210,104)=2^3\cdot3\cdot5\cdot7\cdot13=10920.
\]
Дополнительные множители:
\[
\frac{10920}{210}=52,\qquad \frac{10920}{104}=105.
\]
Тогда
\[
\frac{1}{210}+\frac{3}{104}
=\frac{52}{10920}+\frac{315}{10920}
=\boxed{\frac{367}{10920}}.
\]

\important{Если знаменатели большие, разложение на простые множители помогает одновременно найти НОК и увидеть дополнительные множители.}

\subsection{Сокращение дроби через НОД}
Чтобы сократить дробь максимально, найдите НОД числителя и знаменателя и разделите оба числа на него.

Например,
\[
\nodtext(105,84)=21,
\qquad
\frac{105}{84}=\frac{105:21}{84:21}=\boxed{\frac54}.
\]

Если числа уже разложены на простые множители, общие множители можно просто сократить в разложениях.

\subsection{Десятичная дробь $\rightarrow$ обыкновенная}
Количество цифр после запятой задаёт знаменатель:
\[
0{,}38=\frac{38}{100}=\boxed{\frac{19}{50}}.
\]
После перевода обязательно проверьте, можно ли сократить дробь.

\section{6. Что особенно легко перепутать}

\begin{itemize}[itemsep=0.12em,topsep=0.15em]
  \item $1$ --- ни простое, ни составное число.
  \item Признак делимости на $4$: проверяем две последние цифры как одно число.
  \item Для делимости на $6$ нужны одновременно признаки на $2$ и на $3$.
  \item Для делимости на $12$ нужны одновременно признаки на $3$ и на $4$.
  \item Для делимости на $18$ нужны одновременно признаки на $2$ и на $9$.
  \item В формуле количества делителей работаем с показателями степеней.
  \item В алгоритме Евклида НОД --- последний ненулевой остаток.
  \item Для общего знаменателя ищем НОК знаменателей; для максимального сокращения --- НОД числителя и знаменателя.
\end{itemize}

\section{7. Финальная самопроверка}

Перед контрольной убедитесь, что каждый пункт можно объяснить своими словами без подсказки:

{\small
\begin{multicols}{2}
\begin{itemize}[leftmargin=1.7em,itemsep=0.18em,topsep=0.18em]
  \checkitem{Я различаю делитель и кратное.}
  \checkitem{Я помню признаки делимости на $2,3,4,5,6,9,12$ и $18$.}
  \checkitem{Я умею выписывать делители числа без пропусков.}
  \checkitem{Я отличаю простые, составные числа и число $1$.}
  \checkitem{Я раскладываю число на простые множители.}
  \checkitem{Я нахожу количество делителей по степеням в разложении.}
  \checkitem{Я определяю, являются ли два числа взаимно простыми.}
  \checkitem{Я нахожу НОД и НОК по разложению на простые множители.}
  \checkitem{Я применяю алгоритм Евклида для больших чисел.}
  \checkitem{Я определяю чётность выражения без лишних вычислений.}
  \checkitem{Я использую НОК для общего знаменателя и НОД для сокращения дроби.}
\end{itemize}
\end{multicols}
}

\end{document}
